\documentclass[11pt,a4paper]{article}
\usepackage[margin=23mm]{geometry}
\usepackage{amsmath,amssymb,booktabs}
\usepackage[T1]{fontenc}
\usepackage[hidelinks]{hyperref}
\setlength{\parindent}{0pt}
\setlength{\parskip}{6pt}
\begin{document}
\begin{center}
{\Large VSL-B2: joined service and receiver-state closure}\\[4pt]
13 October 2026 --- conditional reference calculations, not an adopted task
\end{center}

\section*{1. What the local receiving/merging relationship specifies}
Use one identified transfer opportunity and one fixed native CTM interval.
Let $x$ be occupancy of the adjacent receiving cell, $H$ its holding capacity,
$Q$ its per-interval capacity, and $\delta=W/V_0$. Occupancies and transfers
are continuous vehicle-equivalent mass; $Q$ has the same per-interval units.
The source receiving ingredient is
\[
 r(x)=\min\{Q,\delta(H-x)\},\qquad 0\leq x\leq H.
\]
It is space in the destination cell that limits acceptance. For available own-
lane upstream traffic $b$, own-lane priority first gives
\[
 u=\min\{b,Q,r(x)\}.
\]
For a locally available mandatory request $w$, the receiver part of the
lane-change allocation is
\[
 g=\min\left\{w,\frac{r(x)-u}{\gamma}\right\},\qquad \gamma\geq1,
 \qquad u+\gamma g\leq r(x).
\]
Here $\gamma$ denotes the early gap-work weight. It corresponds to the source's
gap ratio for this diagnostic; it is not identified with S13's separate late
work parameter. Since $u\leq r$, the residual is nonnegative without an added
clipping convention.

Accepted mass is $g$, not $\gamma g$: the donor loses $g$ and the target gains
$g$. If $o$ is the target's outflow during the interval and no other inflow is
present, its stock balance is
\[
 I=u+g,\qquad x^+=x+I-o,\qquad r^+=r(x^+).
\]
Thus completed merging can affect subsequent receiving through occupancy.
These equations specify a local relationship conditional on supplied motion,
available traffic and outflow. They do not close an autonomous common-VSL task.

The native algorithm computes residual donor motion after lane-change
allocation. A fixed-geometry task must still qualify its available request and
total admissible donor outflow. If its motion contract supplies a shared
sending budget $S_E$, lateral $g$ and remaining own-lane advance $z$ must meet
$g+z\leq S_E$. This is a compatibility requirement, not an unapproved change
to the source algorithm. Temporary residual calculations and final conservation
must count the accepted transfer once.

\vfill
Source context: Carey--Balijepalli--Watling, Sections 2--5, Algorithm A1 and
the shared residual steps of A2. Selected primary text was checked; complete
algorithm adoption, variable-speed transport and a lane-drop boundary are not
certified.\\
\href{https://eprints.whiterose.ac.uk/id/eprint/77315/7/Carey\%20et\%20al\%202013_with_coversheet.pdf}{Institutional author manuscript}.

\newpage
\section*{2. Deriving the original P1 service example from occupancy}
For this reference calculation only, take
\[
 H=30,\quad Q=10,\quad\delta=\tfrac12,\quad b=4,\quad w=6,\quad\gamma=2.
\]
The two receiver occupancies $x=10$ and $x=22$ give
\[
 r(10)=\min\{10,10\}=10,\qquad
 r(22)=\min\{10,4\}=4.
\]
Applying the priority allocation yields
\[
 (u,g)=(4,3)\quad\hbox{and}\quad (u,g)=(4,0).
\]
Six units of available ending-lane mass therefore leave three or six units
still needing to merge, respectively. The receiving-work tests are
$4+2(3)=10$ and $4+2(0)=4$. Only seven or four units of actual mass enter
the target cell.

The comparison can preserve combined continuing mass by placing 22 or 10 units,
respectively, in a farther continuing region. Including the common own-lane
upstream mass, the compared continuing total is
\[
 C=4+10+22=4+22+10=36.
\]
This bookkeeping embeds the local comparison in equal ending/continuing totals.
It is not a reachable equal-observation-history certificate.

To expose the subsequent stock dependence, condition both cases on the same
feasible target outflow $o=2$. The respective next occupancies and receiving
budgets are
\[
 x_1^+=10+4+3-2=15,\quad r_1^+=\tfrac{15}{2};\qquad
 x_2^+=22+4+0-2=24,\quad r_2^+=3.
\]
An outflow of two is compatible with the local CTM context if the next
receiving constraint is two; it is not a proposed fixed drainage law for
the VSL task.

\begin{center}
\begin{tabular}{lrrrrrr}
\toprule
Case & $x$ & $r$ & $u$ & $g$ & $x^+$ & $r^+$\\
\midrule
More receiving space & 10 & 10 & 4 & 3 & 15 & $15/2$\\
Less receiving space & 22 & 4 & 4 & 0 & 24 & 3\\
\bottomrule
\end{tabular}
\end{center}

\textbf{Traffic conclusion.} Occupancy can generate the ten/four receiving
difference in the approved P1 illustration. Immediate own-lane priority does
not make preparation interaction-free: accepted traffic contributes to later
occupancy. This is one balance calculation, not a traffic rollout, a delay
estimate, an optimal-action result or calibrated parameter choice.

\newpage
\section*{3. A current receiving scalar need not predict future receiving}
Now take another admissible reference setting,
\[
 H=40,\quad Q=10,\quad\delta=\tfrac12,
 \quad b=4,\quad w=6,\quad\gamma=2,\quad o=2.
\]
Compare receiver occupancies $x_1=10$ and $x_2=20$. Both currently give
\[
 r_1=r_2=10,\qquad u_1=u_2=4,\qquad g_1=g_2=3,
 \qquad I_1=I_2=7.
\]
Choose farther continuing masses $f_1=20$ and $f_2=10$. The current grouped
inputs are identical: ending availability is six and
\[
 C_1=4+10+20=34=C_2=4+20+10.
\]
After the \emph{same} inflow and outflow,
\[
 x_1^+=10+7-2=15,\quad r_1^+=10;\qquad
 x_2^+=20+7-2=25,\quad r_2^+=\tfrac{15}{2}.
\]
Hence a proposed update using only the common current grouped masses,
current receiving scalar and even the known local inflow/outflow cannot
return both correct next receiving values. The loss arises on the capacity
plateau: $r=Q$ does not identify how much occupied stock lies behind that
same current receiving value.

\textbf{Scope.} This disproves general closure by a current receiving scalar
for the stated stock-based family. It does not rule out another sufficient
scalar, a more restrictive regime, a compact stock augmentation, or a predictor
with informative history. It does not prove full profiles, discrete order,
or a different best common-speed action.

On the strictly space-limited branch, known parameters allow inversion:
\[
 r<Q\quad\Longrightarrow\quad x=H-r/\delta.
\]
That restricted representation still requires autonomous predictions of own
traffic, available merging and outflow, and must remain valid over its whole
advice range. Carrying occupancy explicitly removes this particular plateau
ambiguity; it does not by itself close downstream supply or physical motion.

\textbf{Decision consequence.} The next shared methodological choice concerns
a compact autonomous receiver/motion closure versus a small spatially resolved
local closure. Neither is selected by this proof.

\newpage
\section*{4. Why speed and complete time remain joint qualifications}
The native source chooses cell length $\ell$ and time interval $\Delta t$ so
that
\[
 \ell=V_0\Delta t.
\]
Keeping both physical geometry and numerical interval fixed while changing
advice to a realized free speed $v$ gives
\[
 \frac{v\Delta t}{\ell}=\frac{v}{V_0}.
\]
An uncongested traversal of one cell per interval is therefore not an unchanged
physical-speed contract for every $v$. This relation identifies a transport
qualification, not a selected variable-speed CTM, a capacity-versus-speed law
or permission to resize cells. A finite fixed-geometry model must explain
its available traffic, receiving and outflow under the same action.

For a fixed gap-work weight, local sensitivity is already branch-dependent.
For example, away from switching boundaries,
\[
 g=w\ \Longrightarrow\ \frac{dg}{dv}=\frac{dw}{dv};\qquad
 g=\frac{r-u}{\gamma}\ \Longrightarrow\
 \frac{dg}{dv}=\frac1\gamma\left(\frac{dr}{dv}-\frac{du}{dv}\right).
\]
The derivatives on the right depend on the missing state/motion closure.
When $u=r$ throughout a neighborhood, $g=0$ there. No universal preparation
benefit follows from extra elapsed time.

For one finite cohort, the retained complete-population objective can be
written independently of a particular local model as
\[
 F(v)=\int_0^{T_{\mathrm{clear}}(v)} N_{\mathrm{unexited}}(t;v)\,dt,
 \qquad \overline T(v)=F(v)/M.
\]
$N_{\mathrm{unexited}}$ counts all cohort mass that has not exited, including
travel and waiting; $M$ is fixed. Finite complete clearance must still be
established for the admitted domain. A common travel duration followed by
half a proportional batch clearance reproduces this expression only under
S13's special timing and service conditions.

\textbf{Present disposition.} Local state-dependent receiving and the R1
receiver allocation are inspectable. An autonomous receiving/motion contract,
joined arrival/departure treatment, lane-end completion, MP6 regime and
physical-action derivative are not admitted. S13's conditional certificate
remains unchanged and is not transferred to this candidate.
\end{document}
